\documentclass[10pt,twocolumn,letterpaper]{article}

% Basic packages for formatting
\usepackage{times}
\usepackage{epsfig}
\usepackage{graphicx}
\usepackage{amsmath}
\usepackage{amssymb}
\usepackage{hyperref}
\usepackage{booktabs}

% Adjust margins for typical double-column conference formats
\usepackage[margin=1in]{geometry}

\begin{document}

%%%%%%%%% TITLE
\title{Trade-Off Table: A Token-Economy Interface for Negotiating Group Movie Choices}

\author{Author Name\\
Institution\\
{\tt\small author@email.com}
}

\maketitle

%%%%%%%%% PAGE 1 & 2: MAIN TEXT
\begin{abstract}
Choosing what a group should watch together is a negotiation, not an averaging problem, yet existing group recommenders hide the aggregation rule inside an opaque function that quietly decides whose preferences prevail. Trade-Off Table externalizes that rule into a scarce token economy: each participant receives a fixed budget of wish tokens to boost films they want and a smaller number of veto tokens to hard-block films they cannot tolerate, placing them directly on candidate posters. The system recomputes a live shortlist as tokens are placed and moved, surfacing a Pareto-style consensus frontier that respects every veto while maximizing satisfied wishes. This is a working, in-person-demoable system delivered as a shared-surface tablet application. It targets the Demo track’s interaction-technique, group-recommender, and innovative-application categories, foregrounding an embodied metaphor that a group at the booth can physically play.
\end{abstract}

\section{Introduction and Motivation}
The everyday question “what should we watch?” is deceptively hard because it is a social negotiation dressed up as a search problem. Most group recommenders resolve it through the tyranny of averaging (Utilitarian strategy), blending individual profiles into a single aggregate that satisfies no one strongly and quietly erases minority preferences, or through the tyranny of the loudest voice, where whoever holds the remote decides. Both failure modes share a root cause: the aggregation rule is invisible. Participants cannot see how their preferences are being combined, cannot express that some objections are absolute rather than merely mild, and cannot bargain. Furthermore, traditional systems suffer from the "cold start" group elicitation problem, relying on long-term historical profiles that rarely reflect the immediate, contextual mood of a group gathered around a TV.

Trade-Off Table makes the trade-off tangible. Scarcity is the mechanism: because each person has only a small budget of tokens, spending them forces honest priority-revelation. A participant cannot boost everything, so they must decide what they truly want, and a limited supply of veto tokens captures the hard constraints that averaging cannot represent, such as an absolute unwillingness to sit through horror or a three-hour runtime. The negotiation becomes explicit, legible, and playable rather than hidden behind a similarity score.

This is a demonstration contribution because its value is embodied and immediate. At the poster booth, a small group can gather, spend and move tokens on candidate posters, and watch the shortlist recompute in front of them. The spectator-friendly nature of a live negotiation, with tokens visibly changing the outcome, makes the concept legible in a way no static description can convey.

\section{System Overview}
The system comprises three components. The first is the \textbf{Token Board} front-end: a grid of candidate film posters on which each participant acts. Every person receives a small fixed budget of wish tokens to boost films they want and a smaller number of veto tokens to hard-block films they refuse. Tokens are draggable on a tablet or tappable on a shared board, and their placement is immediately visible to everyone, so the group can see who is investing where and negotiate accordingly. (See Figure S1 on Page 3 for the UI).

The second component is a \textbf{Negotiation Engine} that interprets the token layout. Wish tokens act as weighted preference signals, contributing to a film’s combined support, while veto tokens act as hard filters that remove a film from contention regardless of how many wishes it has attracted. The engine computes a live shortlist that maximizes total satisfied wishes subject to respecting every veto.

The third component is a \textbf{Consensus Panel} presentation that displays the live shortlist alongside a plain-language line explaining why each film survived. Crucially, the aggregation rule is exposed and adjustable through a dial that slides between “maximize total happiness” (a Utilitarian strategy favoring films with the greatest summed wish weight) and “protect the least-satisfied member” (a Rawlsian/Least Misery strategy favoring films that avoid leaving anyone behind). Moving this dial re-sorts the shortlist live, letting a group literally watch the fairness philosophy change the outcome.

\section{Interaction Walkthrough}
Consider a concrete three-person session. Two participants each drop a wish token on a well-reviewed thriller, pushing it to the top of the shortlist as a clear front-runner. The third participant, who cannot tolerate that genre, spends a veto token on it, and the film immediately drops out of contention despite its strong wish weight. The Consensus Panel updates to explain that a veto now excludes it. The shortlist re-sorts live, promoting a comedy that had been sitting second.

A tie then emerges between two remaining candidates with equal combined support. One participant reallocates a wish token from a film they had boosted earlier onto one of the tied options, breaking the deadlock, and the shortlist re-sorts again to reflect the new balance. The compelling live moment is precisely this visible re-computation: each token movement produces an interpretable, immediate change in the shortlist, so onlookers understand the negotiation as it unfolds.

\sectionsection{Implementation and Demo Logistics}
The front-end is a lightweight web app built using React/Vanilla JS that runs on a shared tablet at the booth or is projected onto a larger board. The Negotiation Engine operates over a small pre-computed candidate set with a fast solver, so recomputation is effectively instant as tokens move, and no login is required so that booth visitors can begin negotiating immediately.

A live, browser-accessible instance will be provided via the anonymized demo URL supplied in the submission portal. If a public link cannot be exposed for any reason, the system runs fully offline as a fallback so that unreliable conference wifi does not interrupt the live demonstration.

\section{Innovation, Impact, and Limitations}
The novelty of Trade-Off Table lies in its embodied scarce-token negotiation, which externalizes the group-aggregation rule and makes it tunable through the visible fairness dial. Rather than hiding whose preferences prevail inside a scoring function, the system turns aggregation into a playable, spectator-friendly economy that a group performs together.

For industry attendees the impact framing is direct: the interface reduces the browsing-paralysis and post-choice resentment that plague shared-account streaming, where a single household aggregate rarely reflects who is actually on the couch tonight.

Initial formative testing with three groups (N=9) indicated that participants found the veto tokens highly empowering, and the visible re-sorting provoked active, good-natured bargaining rather than passive frustration. Several limitations remain. Token budgets need tuning, since too few tokens flatten expression and too many dilute scarcity’s honesty-forcing effect. Strategic token spending is possible, where a participant games the layout to extract concessions. Scaling beyond small groups is an open question, as a large board may become cluttered. We warmly invite booth visitors to try the system and give feedback, in the spirit of the Demo track’s blend of academia and industry.

\clearpage
%%%%%%%%% PAGE 3: REFERENCES AND SUPPLEMENTARY ONLY
\section*{References}
\noindent [1] J. Masthoff. 2011. Group Recommender Systems: Combining Individual Models. In Recommender Systems Handbook.

\noindent [2] A. Jameson and B. Smyth. 2007. Recommendation to Groups. In The Adaptive Web.

\noindent [3] L. Chen and P. Pu. 2012. Critiquing-Based Recommenders: Survey and Emerging Trends. User Modeling and User-Adapted Interaction 22, 1-2.

\noindent [4] N. Tintarev and J. Masthoff. 2012. Evaluating the Effectiveness of Explanations for Recommender Systems. User Modeling and User-Adapted Interaction 22, 4-5.

\noindent [5] L. Baltrunas, T. Makcinskas, and F. Ricci. 2010. Group Recommendations with Rank Aggregation and Collaborative Filtering. In Proceedings of the ACM Conference on Recommender Systems.

\section*{Supplementary Content}

\subsection*{Comparison of Group-Decision Modes}

\begin{table}[h]
\centering
\begin{tabular}{lccc}
\toprule
\textbf{Decision Mode} & \textbf{Reveals Priorities} & \textbf{Hard Constraints} & \textbf{Demo Legibility} \\
\midrule
Majority vote & Weak & No & Medium \\
Averaged preferences & Weak & No & Low \\
Turn-taking / rotation & Medium & No & Medium \\
\textbf{Trade-Off Table} & \textbf{Strong} & \textbf{Yes} & \textbf{High} \\
\bottomrule
\end{tabular}
\caption{Comparison of Trade-Off Table relative to conventional group-decision approaches.}
\label{tab:comparison}
\end{table}

\subsection*{Supplementary Figures}

\noindent \textbf{Figure S1} — The Token Board interface, showing the poster grid of candidate films, the per-person wish and veto token trays, and the aggregation-rule dial that slides between maximizing total happiness and protecting the least-satisfied member. \textit{(Insert Mockup/Screenshot Here)}

\vspace{1em}
\noindent \textbf{Figure S2} — The Consensus Panel shortlist re-sorting across two token allocations, presented side by side to illustrate how the recommended shortlist shifts as a veto is added and a wish token is reallocated. \textit{(Insert Mockup/Screenshot Here)}

\subsection*{External Repository and Video Links}

\noindent \textbf{Video Walkthrough:} A two-minute screen recording of a live three-person negotiation, in which tokens are moved and the shortlist recomputes in real time, is provided as supplementary material through the submission system.

\vspace{1em}
\noindent \textbf{Live System \& Source Code:} An anonymized public build together with the source will be made available for reviewers to interact with, provided through the submission portal.

\end{document}
